\documentclass[10pt,a4paper]{amsart}
\usepackage[margin=19mm]{geometry}
\usepackage{amsmath,amssymb,mathtools}
\usepackage[scr=rsfs,cal=euler]{mathalfa}
\usepackage{xfrac}
\usepackage[unicode,hidelinks,bookmarks=false]{hyperref}
\newcommand{\ie}{i.e.\ }
\newcommand{\cf}{cf.\ }
\newcommand{\resp}{respectively\ }
\newcommand{\Subsection}{\subsection}
\providecommand{\nobreakdash}{\nobreak\hbox{-}}
\setlength{\emergencystretch}{2em}
\multlinegap0pt
\usepackage{siunitx}
\usepackage{siunitx}
\usepackage{siunitx}
\usepackage{siunitx}
\usepackage{tikz}
\usetikzlibrary{cd}
\usepackage{amssymb}
\usepackage{mathtools}
\usepackage{bm}
\usepackage{xcolor}
\usepackage[shortlabels]{enumitem}
\usepackage{float}
\usepackage{tikz}
\usepackage{cancel}
\usepackage{siunitx}
\newtheorem{lemma}{Lemma}
\def\A{\mathcal{A}}
\def\AA{\mathbb{A}}
\def\Gm{\mathbb{G}_\mathrm{m}}
\def\Rgeq{\mathbb{R}_{\geq 0}}
\newcommand{\cal}[1]{\mathcal{#1}}
\def\gridexp{\mathcal{E}_n(X|D)}
\def\logProd{(X|D)^{[n]}}
\newcommand{\logProdexample}[3]{(#1|#2)^{[#3]}}
\newcommand{\pullback}[1]{\arrow[#1, phantom, "\square" , color=black]}
\title{Logarithmic Fulton--MacPherson configuration spaces}
\author{Siao Chi Mok}
\date{Source: arXiv:2503.17563v3 (CC BY 4.0)}
\begin{document}
\maketitle

\noindent\emph{Source and licence.} Logarithmic Fulton--MacPherson configuration spaces. Version arXiv:2503.17563v3 (CC BY 4.0), distributed by arXiv under the Creative Commons Attribution 4.0 International licence, http://creativecommons.org/licenses/by/4.0/, as recorded in the arXiv licence metadata for this exact version and shown on the abstract page https://arxiv.org/abs/2503.17563v3. Full licence notice: \texttt{licenses/arxiv-2503.17563-CC-BY-4.0.txt}.\par\smallskip

\noindent\textit{Editorial scope.} Counted unit: the article's lemma lem:prod\_of\_smooth\_pairs (source0.tex lines 1199--1207). The article's complete original proof of it is delivered with it (source0.tex lines 1208--1238, one \texttt{proof} environment of 31 lines). No mathematics is written, repaired or re-derived: the statement and the proof are the article's own lines, in the article's own notation, and no retained line is substituted or re-flowed.  No further source-local context is needed: every cross-reference of every retained line resolves inside the two delivered spans.  Nothing was omitted from any retained span, so the package carries no omission record.  No retained line contains a citation, so the wrapper carries no bibliography and no bibliography entry.  The retained lines contain the article's own diagram code, which the wrapper typesets with the standard tikz packages; no image file is delivered and the package declares no asset dependency.  Editorial formatting only, in addition: an AMS \texttt{amsart} preamble that restates the article's own declarations and macros used by the retained lines, namely the environments \texttt{lemma} on separate counters, together with the standard packages those lines need.  The retained environments print in this document as Lemma 1 in order of appearance, whereas the article prints the same items with the numbering of its own full document.  The complete native source of the article as distributed in the arXiv e-print for this exact version is bundled unchanged as \texttt{sources/175191/source0.tex}.
\begin{lemma} \label{lem:prod_of_smooth_pairs}
	There exists a Cartesian diagram in the category of schemes and in the fine and saturated logarithmic category
	\begin{equation}\label{cd:prod_of_smooth}
		\begin{tikzcd}
			\logProd \arrow[r] \arrow[d] \pullback{rd} & \prod_{i=1}^{r} \logProdexample{X}{D_i}{n} \arrow[d]\\
			X^n \arrow[r, "\Delta", hook] & (X^n)^r.
		\end{tikzcd}
	\end{equation}
\end{lemma}

\begin{proof}
	Let $\Sigma_i = \Rgeq$ be the tropicalisation of the smooth pair $(X,D_i)$, with Artin fan $\A_i = [\AA^1 / \Gm]$. Observe that we have open embeddings
	\begin{equation*}
		\begin{tikzcd}
			\gridexp \arrow[r,hook] \arrow[d] \pullback{rd} &\prod_{i=1}^{r}\cal{E}_n(X|D_i) \arrow[d]\\
			\A_X^n \arrow[r, hook] & \prod_{i=1}^{r} \A_i^n.
		\end{tikzcd}
	\end{equation*}
%	 which gives the Cartesian diagram
%	\begin{equation*}
%		\begin{tikzcd}
%			\prod_{i=1}^{r}\logProdexample{X}{D_i}{n} \arrow[r] \arrow[d] \pullback{rd} & \gridexp \arrow[d]\\
%			(X^n)^r \arrow[r] & \prod_{i=1}^{r} \A_i^n.
%		\end{tikzcd}
%	\end{equation*}
%	There exists a commutative diagram
%	\begin{equation*}
%		\begin{tikzcd}
%			X^n \arrow[r, "\Delta", hook] \arrow[d] & (X^n)^r \arrow[d]\\
%			\A_X^n \arrow[r, hook] & \prod_{i=1}^{r} \A_i,
% 		\end{tikzcd}
%	\end{equation*} giving a commutative diagram
	The strict morphism $X^n \to \A_X^n \hookrightarrow\prod_{i=1}^{r} \A_i^n$ factors through $X^n \xhookrightarrow{\Delta} (X^n)^r$, giving a commutative diagram
	\begin{equation*}
		\begin{tikzcd}
			\logProd \arrow[r] \arrow[d] & \prod_{i=1}^{r}\logProdexample{X}{D_i}{n} \arrow[r] \arrow[d] \pullback{rd} & \prod_{i=1}^{r}\cal{E}_n(X|D_i) \arrow[d]\\
			X^n \arrow[r, "\Delta", hook] & (X^n)^r \arrow[r, "\text{strict}"] & \prod_{i=1}^{r} \A_i^n
		\end{tikzcd}
	\end{equation*}
	where both the outer and the right squares are Cartesian, thus we obtain the desired Cartesian diagram.
\end{proof}

\end{document}
